\documentclass[12pt]{article}
\usepackage[utf8]{inputenc}
\usepackage[T1]{fontenc}
\usepackage{geometry}
\usepackage{booktabs}
\usepackage{graphicx}
\usepackage{microtype}
\usepackage{parskip}
\usepackage{amsmath}
\usepackage[colorlinks=true, linkcolor=blue, citecolor=blue, urlcolor=blue]{hyperref}
\usepackage[style=ieee, backend=biber]{biblatex}
\addbibresource{references.bib}

\geometry{letterpaper, margin=1in}

\title{CSE 490/590: Project 1 \\[0.5em]
       \large Design and Implementation of a 16-bit Single-Cycle Processor}
\author{
    \small
    Justin Downer \texttt{(jrdowner)} \quad
    Sajjad Ali \texttt{(sali22)} \quad
    Vishal Kishore Kumar \texttt{(vishalki)}
}
\date{April 4, 2026}

\begin{document}
\maketitle

% ---------------------------------------------------------------
\section{Synopsis}
% ---------------------------------------------------------------
This report documents the design, simulation, and hardware synthesis of a 16-bit,
single-cycle, non-pipelined processor \cite{patterson2020}[cite: 1]. The architecture is
byte-addressed with a 16-bit word size and supports 10 mandatory instructions across
R-type, I-type, and J-type formats \cite{patterson2020}[cite: 2]. The processor was developed
using Verilog HDL and verified using both Vivado simulations and the Xilinx Basys~3
FPGA board \cite{digilent2026, xilinx2026}[cite: 3].

% ---------------------------------------------------------------
\section{Work Distribution}
% ---------------------------------------------------------------
While specific modules were contributed by team members, the vast majority of the
architectural design, system integration, and hardware verification was performed
individually[cite: 4].

\begin{table}[h]
  \centering
  \caption{Division of tasks among group members.}
  \label{tab:work-distribution}
  \begin{tabular}{@{}l l p{7.8cm}@{}}
    \toprule
    \textbf{Member} & \textbf{Role} & \textbf{Contribution Details} \\
    \midrule
    Justin Downer & Lead Architect
      & Authored behavioral Verilog for the ALU, Control Unit, and Register
        File \cite{patterson2020}[cite: 5, 6]. \\[4pt]
    Justin Downer & System Integration
      & Developed the structural top-level wrapper using Vivado Block Designer
        and managed all IP interconnections \cite{xilinx2026}[cite: 7]. \\[4pt]
    Justin Downer & Hardware Lead
      & Managed XDC mapping, bitstream generation, and 7-segment display
        integration for hardware demo \cite{digilent2026}[cite: 8]. \\[4pt]
    Team Members  & Contributors
      & Developed supplementary logic modules[cite: 9]. \\
    \bottomrule
  \end{tabular}
\end{table}

% ---------------------------------------------------------------
\section{Implementation Details}
% ---------------------------------------------------------------
\subsection{Design Methodology}
The processor integration utilized \textbf{Vivado's Block Designer} to visualize the
datapath and component interconnections[cite: 10]. A structural Verilog file was maintained via
the \textbf{Integrated HDL Wrapper} feature, ensuring that any changes to the block
design were automatically reflected in the top-level code \cite{xilinx2026}[cite: 11]. This
reduced manual wiring errors while allowing individual components to be designed using
\textbf{behavioral Verilog} \cite{patterson2020}[cite: 12].

\subsection{Instruction Fetching and Memory Alignment}
To maintain byte-addressing, the \textbf{Program Counter (PC)} increments by~2 for
each sequential instruction \cite{patterson2020}[cite: 13]. An \textbf{Adder} module computes
the next PC value before it is passed to Instruction Memory[cite: 14]. The Instruction Memory is
initialized via the \texttt{readmemh} command, loading \texttt{vbit.mem} and slicing
the 16-bit output into appropriate bit-fields for the Control Unit, Register File, and
Sign Extension \cite{patterson2020}[cite: 15].

\subsection{Branch and Jump Control Logic}
The branching mechanism uses a specialized \textbf{XOR-based logic path} to handle
both \texttt{BEQ} and \texttt{BNE} instructions \cite{patterson2020}[cite: 16]. The \textbf{Zero
Flag} from the ALU is XORed with the \texttt{BNE} signal from the Control Unit[cite: 17];
the result is ANDed with the Branch Flag to drive the \texttt{PCSrc} multiplexer[cite: 18]. A custom
\textbf{ShiftLeft1} module doubles the effective range of signed immediates while
keeping the PC word-aligned \cite{patterson2020}[cite: 19].

\subsection{Sign Extension and ALU Control}
Two distinct \textbf{Sign Extension} modules handle the 4-bit signed immediate
(I-type) and the 12-bit signed immediate (J-type) \cite{patterson2020}[cite: 20]. A dedicated
\textbf{ALU Control} module decodes function bits for R-type instructions
(e.g.,~\texttt{ADD}, \texttt{SUB}, \texttt{SLL}, \texttt{AND}), improving modularity
and simplifying the primary Control Unit \cite{patterson2020}[cite: 21].

% ---------------------------------------------------------------
\section{Simulation Results}
% ---------------------------------------------------------------
The following waveforms demonstrate the functional correctness of the 10 mandatory instructions within the Vivado simulation environment[cite: 22, 23].
\subsection{CPU Block Diagram}
\begin{figure}[ht]
    \centering
    \includegraphics[width=0.95\textwidth]{CPU_diagram.pdf}
    \caption{Top-level structural block diagram generated from Vivado IP Integrator.}
    \label{fig:block_diagram}
\end{figure}

\clearpage
\subsection{Instruction Waveforms}

\begin{figure}[ht]
    \centering
    \includegraphics[width=0.85\textwidth]{ADD_Waveform.png}
    \caption{Waveform for R-type \texttt{ADD} instruction.}
\end{figure}

\begin{figure}[ht]
    \centering
    \includegraphics[width=0.85\textwidth]{SUB_Waveform.png}
    \caption{Waveform for R-type \texttt{SUB} instruction.}
\end{figure}

\begin{figure}[ht]
    \centering
    \includegraphics[width=0.85\textwidth]{SLL_Waveform.png}
    \caption{Waveform for R-type \texttt{SLL} instruction demonstrating shift-left logic.}
\end{figure}

\begin{figure}[ht]
    \centering
    \includegraphics[width=0.85\textwidth]{AND_Waveform.png}
    \caption{Waveform for R-type \texttt{AND} instruction.}
\end{figure}

\begin{figure}[ht]
    \centering
    \includegraphics[width=0.85\textwidth]{LW_Waveform.png}
    \caption{Waveform for I-type \texttt{LW} instruction showing memory-to-register data transfer.}
\end{figure}

\begin{figure}[ht]
    \centering
    \includegraphics[width=0.85\textwidth]{SW_Waveform.png}
    \caption{Waveform for I-type \texttt{SW} instruction showing register-to-memory data transfer.}
\end{figure}

\begin{figure}[ht]
    \centering
    \includegraphics[width=0.85\textwidth]{ADDI_Waveform.png}
    \caption{Waveform for I-type \texttt{ADDI} instruction using the 4-bit sign-extended immediate.}
\end{figure}

\begin{figure}[ht]
    \centering
    \includegraphics[width=0.85\textwidth]{BEQ_Waveform.png}
    \caption{Waveform for I-type \texttt{BEQ} instruction verifying PC-relative branching.}
\end{figure}

\begin{figure}[ht]
    \centering
    \includegraphics[width=0.85\textwidth]{BNE_Waveform.png}
    \caption{Waveform for I-type \texttt{BNE} instruction verifying the XOR-based branch logic.}
\end{figure}

\begin{figure}[ht]
    \centering
    \includegraphics[width=0.85\textwidth]{JMP_Waveform.png}
    \caption{Waveform for J-type \texttt{JMP} instruction demonstrating PC-relative jump addressing.}
\end{figure}

\clearpage
% ---------------------------------------------------------------
\section{Hardware Verification}
% ---------------------------------------------------------------
\subsection{Memory Initialization and Test Strategy}
Data Memory was initialized using a Verilog \texttt{initial} block to hardcode test
values for immediate verification of \texttt{LW} and \texttt{SW} instructions
\cite{patterson2020}[cite: 22]. Success is defined as the architectural state matching expected
values across all 10 mandatory instructions \cite{patterson2020}[cite: 23].

\subsection{Hardware Visualization}
To bridge the gap between internal logic and physical observation, the 16-bit
\textbf{ALUResult} is mapped to the Basys~3 board's \textbf{7-segment display},
providing a hexadecimal representation of calculations during the live demonstration
\cite{digilent2026}[cite: 24].

\subsection{Synthesis and Implementation}
The design passed through the Vivado \textbf{Synthesis} and \textbf{Implementation}
flows without timing violations \cite{xilinx2026}[cite: 25]. The absence of critical errors
confirms that the logic is physically routable on the Artix-7 FPGA \cite{xilinx2026}[cite: 26].

% ---------------------------------------------------------------
\section{Conclusion}
% ---------------------------------------------------------------
The successful implementation of this 16-bit processor demonstrates a complete design
cycle — from ISA specification to FPGA synthesis[cite: 27]. The modular approach ensured stable
execution in both simulation and hardware environments \cite{patterson2020, xilinx2026}[cite: 27].

\printbibliography
\end{document}